\documentclass[lettersize,journal]{IEEEtran}
\usepackage{listings} 
\usepackage{xcolor} 
\usepackage[utf8]{inputenc}  
\usepackage{amsmath,amsfonts}
\usepackage{algorithmic}
\usepackage{array}
\usepackage[caption=false,font=normalsize,labelfont=sf,textfont=sf]{subfig}
\usepackage{fancyvrb}
\usepackage{textcomp}
\usepackage{stfloats}
\usepackage{url}
\usepackage{verbatim}
\usepackage{graphicx}
\usepackage{ctex}
\hyphenation{op-tical net-works semi-conduc-tor IEEE-Xplore}
\def\BibTeX{{\rm B\kern-.05em{\sc i\kern-.025em b}\kern-.08em
		T\kern-.1667em\lower.7ex\hbox{E}\kern-.125emX}}
\usepackage{balance}
\usepackage{texshade}
\usepackage{fontspec}    % For XeLaTeX  
\usepackage{xeCJK}       % For Chinese support  

% 设置中文字体  
\setCJKmainfont{SimSun} 

% 设置代码高亮  
\lstset{  
	language=Python,                % 选择语言  
	basicstyle=\ttfamily\small,    % 基本字体样式  
	keywordstyle=\color{blue},      % 关键字颜色  
	commentstyle=\color{green},     % 注释颜色  
	stringstyle=\color{red},        % 字符串颜色  
	numbers=left,                   % 行号显示在左侧  
	numberstyle=\tiny\color{gray},  % 行号样式  
	stepnumber=1,                   % 每行显示行号  
	numbersep=5pt,                  % 行号与代码之间的距离  
	backgroundcolor=\color{white},  % 背景颜色  
	showspaces=false,               % 不显示空格  
	showstringspaces=false,         % 不显示字符串中的空格  
	showtabs=false,                 % 不显示制表符  
	tabsize=2,                      % 制表符宽度  
	breaklines=true                  % 自动换行  
}  


\begin{document}
	
	
\title{三维空间管道布局课设报告}
\author{\\华中科技大学，人工智能与自动化学院，湖北武汉430074
\thanks{}}
	
	
\maketitle
	
\begin{abstract}
本文章旨在通过强化学习算法解决三维空间管道布局优化问题。通过定义一个三维空间矩阵，其中包含管道和障碍物的位置，我们的目标是找到管道的最佳铺设方式，同时满足特定的管道连接规则
\end{abstract}
	
\begin{IEEEkeywords}
DQN算法；三维空间管道布局；贪心策略
\end{IEEEkeywords}
	
	
\section{引言}
\subsection{功能需求}
根据QT某系统空间管道布局，我们需要在三维空间中连接多条管道，同时满足特定的管道连接规则。

\subsection{数据提供}
已知三维空间被划分为$i \times j \times k$个小立方体，每个立方体标记为0或1，0表示未利用空间，1表示包含零部件。

\subsection{管道连接规则}
管道连接需遵循以下规则：
\begin{enumerate}
	\item 避免碰撞；
	\item 直管段长度要求；
	\item 90°拐弯角度；
	\item 连续弯之间的直线管段长度；
	\item 管道间固定宽度的间距；
	\item 横平竖直的布局；
	\item 最小化弯点数量。
\end{enumerate}


\section{DQN算法的解释}
	\subsection{背景}
	DQN是由DeepMind在2015年提出的，它通过使用深度神经网络来逼近Q值函数，以解决高维状态空间的问题。Q值函数表示在给定状态下采取某一动作的期望回报。
	\subsection{DQN的基础}  
	DQN的核心基于Q-learning，Q-learning的基本思想是：
\begin{itemize}  
	\item \textbf{Q值}：$ Q(s, a) $ 代表在状态 $ s $ 下采取动作 $ a $ 的期望回报。  
	\item \textbf{Bellman方程}：更新Q值的方式如下：  
	\begin{equation}  
		Q(s, a) \leftarrow Q(s, a) + \alpha [r + \gamma \max_{a'} Q(s', a') - Q(s, a)]  
	\end{equation}  
	其中，$ r $ 是即时奖励，$ s' $ 是执行动作后到达的新状态，$ \alpha $ 是学习率，$ \gamma $ 是折扣因子。  
\end{itemize}  

	\subsection{DQN的主要组件}
	\subsubsection{经验回放（Experience Replay）}  
	\begin{itemize}  
		\item \textbf{目的}：通过存储智能体经历的转换（状态、动作、奖励、下一个状态）来打破数据之间的相关性。  
		\item \textbf{实现}：将每次经历存储在缓冲区中，随机抽取小批量样本进行训练。  
	\end{itemize}  
	
	\subsubsection{目标网络（Target Network）}  
	\begin{itemize}  
		\item \textbf{目的}：解决DQN中的不稳定性问题，避免直接使用当前的Q网络进行目标更新。  
		\item \textbf{实现}：维护一个与主网络结构相同但参数不同的网络，每隔一段时间（例如每固定步数）同步一次参数。  
	\end{itemize}  
	
	\subsubsection{ε-贪婪策略（ε-greedy Policy）}  
	\begin{itemize}  
		\item \textbf{目的}：平衡探索（exploration）和利用（exploitation）。  
		\item \textbf{实现}：以概率 $ \epsilon $ 随机选择动作（探索），以概率 $ 1 - \epsilon $ 选择当前Q网络中评分最高的动作（利用）。  
	\end{itemize}  
	
    \subsection{DQN的训练流程}  
    \begin{enumerate}  
	\item \textbf{初始化}：  
	\begin{itemize}  
		\item 初始化主Q网络和目标Q网络，随机初始化网络参数。  
		\item 初始化经验回放缓冲区。  
	\end{itemize}  
	
	\item \textbf{循环训练}：  
	\begin{itemize}  
		\item 在每个时刻 $ t $：  
		\begin{enumerate}  
			\item \textbf{选择动作}：根据当前状态 $ s_t $ 使用ε-贪婪策略选择动作 $ a_t $。  
			\item \textbf{执行动作}：在环境中执行动作，观察即时奖励 $ r_t $ 和下一个状态 $ s_{t+1} $。  
			\item \textbf{存储经验}：将转换 $ (s_t, a_t, r_t, s_{t+1}) $ 存储到经验回放缓冲区。  
			\item \textbf{采样小批量}：从经验回放中随机抽取小批量样本 $ (s, a, r, s') $。  
			\item \textbf{计算目标}：  
			\begin{equation}  
				y = r + \gamma \max_{a'} Q_{target}(s', a')  
			\end{equation}  
			\item \textbf{更新Q网络}：使用随机样本更新主Q网络的参数，最小化损失：  
			\begin{equation}  
				L = (y - Q(s, a; \theta))^2  
			\end{equation}  
			\item \textbf{每隔固定步数更新目标网络}：将当前主网络的权重复制到目标网络。  
		\end{enumerate}  
	\end{itemize}  
	
	\item \textbf{训练结束}：在满意的回报或某个训练周期后停止训练，根据环境的反馈进行评估。  
\end{enumerate}  

	
	\section{代码实现}
	\subsection{运行思路}
	程序主要由五个部分组成：env/pipe\_env.py负责生成三维空间
	
	\subsubsection{三维空间初始化}
	在 \texttt{env/pipe\_env.py} 中定义的 \texttt{PipeLayoutEnv} 类负责生成三维空间：  
	\begin{itemize}  
		\item 0 表示可用空间。  
		\item 1 表示障碍（零部件）。  
		起点（start）和终点（goal）也在此初始化。 
		\item  代码示例：
			\begin{lstlisting}[language=Python]   
			space = np.zeros(config["space_shape"])  
			space[3:5, 3:5, 3:5] = 1  # 添加障碍物  
			env = PipeLayoutEnv(space, config["start"], config["goal"])  
		\end{lstlisting}  
		\item  思路：
		空间由 numpy 数组表示，方便操作。障碍物可以通过直接修改数组值来添加，例如 \texttt{space[x, y, z] = 1}。  
	\end{itemize}  

	\subsubsection{环境模拟 (PipeLayoutEnv 类)}  
	强化学习智能体与环境交互时，每个动作会更新管道路径。
	     
	\texttt{step(action)}:  
	\begin{itemize}  
		\item 接收智能体的动作（如上下左右前后）。  
		\item 检查新位置是否越界或碰撞。  
		\item 更新状态、路径和奖励。  
	\end{itemize}  
	
	\texttt{reset()}:  
	\begin{itemize}  
		\item 每轮训练开始时，重置环境到初始状态。  
		\item 代码示例：
		\begin{lstlisting}  
			# 示例: 执行动作  
			next_state, reward, done = env.step(action)  
		\end{lstlisting}  
		\item 思路：
		环境充当智能体的评估工具，计算动作的有效性（是否碰撞、越界等）。通过奖励值引导智能体寻找最优路径。 
	\end{itemize}    
	
	\subsubsection{强化学习训练 (train.py 文件)} 
	
	核心逻辑是DQN（深度Q网络）：  
	\begin{itemize}  
		\item \textbf{policy\_net}: 当前网络，用于选择动作。  
		\item \textbf{target\_net}: 目标网络，用于计算目标Q值（用来稳定训练）。  
	\end{itemize}  
	
    经验回放 (Replay Memory)：  
	\begin{itemize}  
		\item 存储智能体与环境的交互数据（状态、动作、奖励、下一状态）。  
		\item 从中随机抽样用于训练，打破数据相关性。  
		\item 代码示例 :
		\begin{lstlisting}[language=Python]   
			# 每轮训练过程  
			state = env.reset()  
			for step in range(max_steps):  
			action_idx = select_action(state, epsilon)  # 选择动作  
			next_state, reward, done = env.step(actions[action_idx])  # 与环境交互  
			memory.append((state, action_idx, reward, next_state, done))  # 存储经验  
			state = next_state  
			if done:  
			break  
		\end{lstlisting}  
			\item 思路：
			
				智能体通过Q网络评估各动作的价值(Q值)，选择最优动作。
				
				每轮训练后，智能体逐步学会避开障碍、找到终点。  
	\end{itemize}  
	\subsubsection{奖励函数设计}

奖励函数引导智能体优化路径，具体包括：  
\begin{itemize}  
	\item \textbf{获取奖励}：路径经过障碍时，给予负奖励。  
	\item \textbf{达到终点奖励}：路径到达终点时，给予大额度正奖励。  
	\item \textbf{步长惩罚}：每步默认轻微惩罚，鼓励尽快找到终点。  
	\item \textbf{美观性奖励（可选）}：  
	\begin{itemize}  
		\item 减少拐点，保持直线路径。  
	\end{itemize}  
	\item 代码示例
	\begin{lstlisting}[language=Python]   
		if new_pos == self.goal:  
		reward = 100  
		elif collision:  
		reward = -10  
		else:  
		reward = -1  # 无偏好的普通选择  
	\end{lstlisting} 
		\item 思路：
		
		奖励机制是强化学习的核心。设计合理的奖励策略能有效地引导智能体找到最优路径。  
\end{itemize}  

	\subsubsection{动作选择策略}

	动作空间包含 6 个方向（上下左右前后）。初始阶段：  
	\begin{itemize}  
		\item 以一定概率选择随机动作（探索）。  
		\item 则余概率选择最优动作（利用）。  
	\end{itemize}  
	
	随着训练进行：  
	\begin{itemize}  
		\item 随机动作概率逐渐减少（ε-贪心策略），智能体能够更多利用已学到的知识。
		\item 代码示例
		\begin{lstlisting}[language=Python]   
			def select_action(state, epsilon):  
			if random.random() < epsilon:  
			return random.randint(0, 5)  # 随机动作  
			else:  
			state_tensor = torch.tensor(state, dtype=torch.float32)  
			return policy_net(state_tensor).argmax().item()  # 选择动作  
		\end{lstlisting}  
			\item 思路：
			
			初期随机探索，避免局部最优。后期依赖学习到的策略，优化效率。
	\end{itemize}  

	
	\subsubsection{ 可视化路径 (utils/visualize.py 文件)}  
	
	训练完成后，调用 \texttt{visualize\_path()} 函数展示路径。  
	\begin{itemize}
	 \item 代码示例:
	
	 \begin{lstlisting}[language=Python]  
	 path = np.array(env.path)  # 获取路径  
	 fig = plt.figure()  
	 ax = fig.add_subplot(111, projection='3d')  
	 ax.plot(path[:, 0], path[:, 1], label="Pipe Path")  
	 plt.legend()  
	 plt.show()  
	 \end{lstlisting}  
	  \item 思路:
	   
	提取路径点坐标，并用 matplotlib 的三维图形展示结果，直观反映路径走向。  
\end{itemize} 
 
	
	\subsubsection{ 参数调整 (config.py 文件)}  
	通过修改 config.py 中的参数，可以调整以下配置：  
	\begin{itemize}  
		\item \textbf{三维空间大小}：  
		\begin{lstlisting} [language=Python]  
			"space_shape": (10, 10, 10)  # 可设为 (20, 20, 20)  
		\end{lstlisting}  
		\item \textbf{起点和终点}：  
		\begin{lstlisting} [language=Python]  
			"start": (0, 0, 0),  
			"goal": (9, 9, 9),  
		\end{lstlisting}  
		\item \textbf{训练轮数}：  
		\begin{lstlisting} [language=Python]  
			"num_episodes": 1000,  # 训练为500以强化训练  
		\end{lstlisting}  
		\item \textbf{学习率和奖励折扣因素}：  
		\begin{lstlisting} [language=Python]  
			"lr": 0.001,  
			"gamma": 0.99,  
		\end{lstlisting}  
	\end{itemize}  
	
	\subsection{运行方式}
	\begin{enumerate}
		\item 定义网格的宽度和高度，以及障碍物的位置。
		\item 初始化`GridEnvironment`和`AStar`对象。
		\item 生成随机的起点和终点。
		\item 对每对起点和终点，使用`AStar`对象搜索路径，并将结果存储在路径列表中。
		\item 调用`draw\_paths`函数，绘制并显示路径规划结果。
	\end{enumerate}
	
	\subsection{源代码}
	\begin{lstlisting}[language=Python]  
		main.py
		from env.pipe_env import PipeLayoutEnv
		from train import train_dqn
		from utils.visualize import visualize_path
		import numpy as np
		
		if __name__ == "__main__":
		from config import config
		
		# 初始化空间
		space = np.zeros(config["space_shape"])
		space[3:5, 3:5, 3:5] = 1  # 添加障碍物
		env = PipeLayoutEnv(space, config["start"], config["goal"])
		
		# 训练
		memory = []
		train_dqn(env, config["num_episodes"], memory, config["batch_size"])
		
		# 可视化
		visualize_path(env.path)
		
		train.py
		import random
		import torch
		import torch.optim as optim
		import torch.nn.functional as F
		from models.dqn_model import DQN
		from env.pipe_env import PipeLayoutEnv
		
		# 超参数
		gamma = 0.99
		epsilon = 1.0
		epsilon_decay = 0.995
		min_epsilon = 0.01
		lr = 0.001
		
		def train_dqn(env, num_episodes, memory, batch_size):
		policy_net = DQN(3, 6)  # 输入3维坐标，输出6个方向
		target_net = DQN(3, 6)
		target_net.load_state_dict(policy_net.state_dict())
		optimizer = optim.Adam(policy_net.parameters(), lr=lr)
		
		for episode in range(num_episodes):
		state = env.reset()
		total_reward = 0
		done = False
		
		while not done:
		# 选择动作
		if random.random() < epsilon:
		action_idx = random.randint(0, 5)
		else:
		state_tensor = torch.tensor(state, dtype=torch.float32)
		with torch.no_grad():
		action_idx = policy_net(state_tensor).argmax().item()
		
		# 执行动作
		action = [(1, 0, 0), (-1, 0, 0), (0, 1, 0), (0, -1, 0), (0, 0, 1), (0, 0, -1)][action_idx]
		next_state, reward, done = env.step(action)
		
		# 记录经验
		memory.append((state, action_idx, reward, next_state, done))
		if len(memory) > 10000:
		memory.pop(0)
		
		# 经验回放
		if len(memory) >= batch_size:
		batch = random.sample(memory, batch_size)
		states, actions, rewards, next_states, dones = zip(*batch)
		
		states = torch.tensor(states, dtype=torch.float32)
		actions = torch.tensor(actions, dtype=torch.long)
		rewards = torch.tensor(rewards, dtype=torch.float32)
		next_states = torch.tensor(next_states, dtype=torch.float32)
		dones = torch.tensor(dones, dtype=torch.float32)
		
		q_values = policy_net(states).gather(1, actions.unsqueeze(1)).squeeze()
		next_q_values = target_net(next_states).max(1)[0]
		expected_q_values = rewards + gamma * next_q_values * (1 - dones)
		
		loss = F.mse_loss(q_values, expected_q_values)
		optimizer.zero_grad()
		loss.backward()
		optimizer.step()
		
		state = next_state
		total_reward += reward
		
		if episode % 10 == 0:
		target_net.load_state_dict(policy_net.state_dict())
		
		print(f"Episode {episode}: Total Reward = {total_reward}")
		
		pipe_env.py
		import numpy as np
		
		class PipeLayoutEnv:
		def __init__(self, space, start, goal):
		self.space = space  # 三维空间
		self.start = start  # 起点
		self.goal = goal    # 终点
		self.current_pos = start
		self.path = [start]  # 已生成的路径
		
		def step(self, action):
		"""执行动作并返回新的状态、奖励、是否完成"""
		x, y, z = self.current_pos
		dx, dy, dz = action
		new_pos = (x + dx, y + dy, z + dz)
		
		if not (0 <= new_pos[0] < self.space.shape[0] and
		0 <= new_pos[1] < self.space.shape[1] and
		0 <= new_pos[2] < self.space.shape[2]):
		return self.current_pos, -10, False  # 越界惩罚
		
		if self.space[new_pos] == 1:
		return self.current_pos, -10, False  # 碰撞惩罚
		
		self.current_pos = new_pos
		self.path.append(new_pos)
		reward = 1 if new_pos != self.goal else 100
		done = new_pos == self.goal
		return new_pos, reward, done
		
		def reset(self):
		"""重置环境"""
		self.current_pos = self.start
		self.path = [self.start]
		return self.start
		
		dqn_model.py
		import torch
		import torch.nn as nn
		
		class DQN(nn.Module):
		def __init__(self, input_dim, output_dim):
		super(DQN, self).__init__()
		self.fc = nn.Sequential(
		nn.Linear(input_dim, 128),
		nn.ReLU(),
		nn.Linear(128, 128),
		nn.ReLU(),
		nn.Linear(128, output_dim)
		)
		
		def forward(self, x):
		return self.fc(x)
		
		config.py
		config = {
			"space_shape": (10, 10, 10),
			"start": (0, 0, 0),
			"goal": (9, 9, 9),
			"num_episodes": 1000,
			"batch_size": 64,
			"lr": 0.001,
			"gamma": 0.99
		}
		
		visualize.py
		import matplotlib.pyplot as plt
		from mpl_toolkits.mplot3d import Axes3D
		import numpy as np
		
		def visualize_path(path):
		path = np.array(path)
		fig = plt.figure()
		ax = fig.add_subplot(111, projection='3d')
		ax.plot(path[:, 0], path[:, 1], path[:, 2], label='Pipe Path')
		plt.legend()
		plt.show()
	\end{lstlisting} 
	
	
	\begin{figure}[!t]
		\centering
		\includegraphics[width=3in]{fig1}
		\caption{1000轮运行结果}
		\label{}
	\end{figure} 
	
	\begin{figure}[!t]
		\centering
		\includegraphics[width=3in]{fig2}
		\caption{50轮运行结果}
		\label{}
	\end{figure} 
	\section{总结}
该项目实现了一个能够模拟三维空间管道布局的环境\texttt{PipeLayoutEnv}，它能够处理管道的移动、碰撞检测和路径记录。项目设计并训练了一个DQN模型，该模型能够学习在给定的三维空间中找到最优路径。通过经验回放机制，模型能够在与环境的交互中有效地利用历史数据，提高了学习的效率。引入了目标网络来稳定训练过程，减少了模型在训练过程中的波动。可视化了智能体找到的路径，直观地展示了模型的性能。对模型的性能进行了评估，包括路径长度、碰撞次数和路径的美观性。


	\begin{thebibliography}{1}
		\bibliographystyle{IEEEtran}
		
		\bibitem{ref1}
    	V. Mnih et al., "Human-level control through deep reinforcement learning," Nature, vol. 518, no. 7540, pp. 529-533, 2015/02/01 2015.
		
		\bibitem{ref2}
		V. Mnih et al., "Playing Atari with Deep Reinforcement Learning," p. arXiv:1312.5602Accessed on: December 01, 2013. doi: 10.48550/arXiv.1312.5602 Available: https://ui.adsabs.harvard.edu/abs/2013arXiv1312.5602M		
		
	\end{thebibliography}
\end{document}
